% Below Macros for diagram generation generated by Claude on 2026-10-18
% (reworked on 2026-10-20 to use TikZ pics and tikz-3dplot; number lines,
% vectors, parallelograms, planes and parallelepipeds added 2026-10-20)

% ------------------------------------------------------------
% Figures for linear algebra
% (in the style of Lay, Linear Algebra and Its Applications)
%
% Built on standard TikZ machinery:
%   - each figure is a TikZ pic, so several can share one tikzpicture
%     (with arrows, labels, etc. drawn between them);
%   - 3D figures use tikz-3dplot's main coordinate system for the view;
%   - options are ordinary pgfkeys.
% Requires \usepackage{tikz}; loads tikz-3dplot and the calc library
% itself. \input this file in the preamble. Examples: diagram_examples.tex.
%
% Pics (use inside a tikzpicture):
%   2D
%   \pic {number line={keys}};     number line with points and intervals
%   \pic {axes 2d={keys}};         x1x2-axes, optional grid and tick labels
%   \pic {vector={keys}};          arrow from `from' to `to', label at tip
%   \pic {parallelogram={keys}};   region spanned by u and v
%   \pic {basis grid={keys}};      grid lines along u and v (rotated, skewed,
%                                  or the image of the standard grid under A)
%   \pic {ellipse={keys}};         ellipse x1^2/a^2 + x2^2/b^2 <= 1
%   3D (put llfig view on the tikzpicture so they share one view)
%   \pic {axes 3d={keys}};         x1x2x3-axes
%   \pic {vector={keys}};          as in 2D, with 3 coordinates
%   \pic {parallelogram={keys}};   as in 2D, with 3 coordinates
%   \pic {basis grid={keys}};      as in 2D: a grid on the plane of u and v
%   \pic {plane={keys}};           patch of the plane spanned by u and v
%   \pic {parallelepiped={keys}};  solid spanned by u, v, w (hidden edges dashed)
%   \pic {ellipsoid={keys}};       x1^2/a^2 + x2^2/b^2 + x3^2/c^2 <= 1
% \pic at (p) {...} places a pic with its origin at p, e.g. a plane
% through p.
%
% 3D view, for the whole tikzpicture:
%   \begin{tikzpicture}[llfig view]             default view (30, 20)
%   \begin{tikzpicture}[llfig view={45}{30}]    {azimuth}{elevation}
% x1 points toward the viewer, `azimuth' degrees left of straight on; the
% eye is `elevation' degrees up. Same as tikz-3dplot's
% \tdplotsetmaincoords{90-elevation}{90+azimuth} + tdplot_main_coords.
%
% Enlarging a picture that contains pics: use llfig scale, not scale.
%   \begin{tikzpicture}[llfig scale=1.5]
%   \begin{tikzpicture}[llfig view, llfig scale=1.5]   (after llfig view)
% Like nodes, TikZ pics ignore a picture's scale=..., so pics and plain
% \draw commands would drift apart. llfig scale multiplies the coordinate
% units instead, which pics do follow; text keeps its size.
%
% Linear maps: transform everything drawn in a scope by a matrix A,
% given by rows. Coordinates x land at A x; grids, axes, vectors, pics,
% curves and circles (which become ellipses) all follow, while text,
% dots, arrowheads and line widths keep their shape.
%   \begin{scope}[llfig matrix={2,-1}{1,1}] ... \end{scope}        2x2
%   \begin{scope}[llfig matrix 3d={1,0,0.5}{0,1,0}{0,0,1}] ... \end{scope}
% Entries may be expressions, e.g. {cos(30),-sin(30)}{sin(30),cos(30)}.
% Nested scopes compose: the inner matrix is applied first. Not
% transformed: the `rectangle' and `grid' path operations (they always
% draw page-aligned shapes; use -- cycle and the basis grid pic instead)
% and the ellipsoid pic.
%   \SubspaceFig[keys]{{1,2},{0,1}}
% draws Span of the given columns in R^1, R^2 or R^3 (the number of
% entries per column). The dimension is computed (rank of the columns),
% so a dependent set is drawn as what it really spans: a point, a line
% through 0, a plane (a square patch from a Gram-Schmidt orthonormal
% basis of the span in R^3; the whole window in R^2), or all of R^3.
% The columns are drawn as arrows v_1, v_2, ... Keys: name (e.g. $W$),
% name pos (where a plane's name goes, as for the plane pic),
% vector name (default v), window (half-width shown, default 2.5),
% color (default blue), dim caption=true|false (caption "dim W = r"
% below; default true; label below overrides it), ambient=1|2|3 (needed
% only for the zero subspace, given as {}), azimuth, elevation, scale.
%
%   \KernelRangeFig[keys]
% Lay-style schematic (Sec. 4.2): blobs for R^n and R^m, Nul A inside the
% first and Col A inside the second, x |-> Ax, and Nul A |-> 0. Keys:
% domain, codomain, kernel, range, map label, point, image (the labels
% x and Ax); e.g. domain=$V$, codomain=$W$, kernel={$\ker T$},
% range={$\operatorname{range} T$}, map label={$T$},
% image={$T(\mathbf{x})$} for a general linear transformation.
%   \FundamentalSubspacesFig[keys]
% Lay-style schematic (Sec. 6.1): Row A with the perpendicular Nul A in
% R^n, Col A with the perpendicular Nul A^T in R^m. Keys: domain,
% codomain, row space, kernel, range, left null, map label, dims=true
% (adds dim r, n - r, r, m - r), rank (the symbol for r, default r).
%   \NulColFig[keys]{{rows of A}}
% The actual Nul A (in R^n) and Col A (in R^m) of a matrix given by rows,
% m, n <= 3, e.g. \NulColFig{{1,0,1},{0,1,1}}. Nul A is computed as
% (Row A)^perp: with row rank 2 in R^3 it is spanned by the cross product
% of two independent rows, with rank 1 by two vectors perpendicular to
% the row. Keys: matrix name (default A), window, azimuth, elevation.
%
%   \MatrixMapFig[keys]{row 1}{row 2}{content}
% draws the TikZ content, an arrow, and the content again under A.
% Keys: label (over the arrow, default $T$), before label, after label
% (captions under the two pictures), scale. Needs amsmath (\xrightarrow).
%
% Stand-alone shortcuts (each is its own tikzpicture, baseline at the
% origin so side-by-side figures line up along the x1-axis):
%   \EllipseFig[keys]     \UnitDiskFig[keys]    (a = b = 1, intercepts labeled 1)
%   \EllipsoidFig[keys]   \UnitSphereFig[keys]  (a = b = c = 1, intercepts labeled 1)
%   \SpheroidFig[keys]    (b = a: equatorial radius a, polar radius c)
%   \NumberLineFig[keys]
%   \ParallelogramFig[keys]    parallelogram on axes sized to fit
%   \ParallelepipedFig[keys]   parallelepiped on 3D axes sized to fit
%
% Keys (all optional; each pic uses the ones that apply to it):
%   General
%   x1 label, x2 label, x3 label
%                           axis labels (default $x_1$, $x_2$, $x_3$; empty = none)
%   name                    region label, e.g. $S$ or $T(S)$ (default none)
%   label above, label below
%                           caption centred over / under the figure's origin,
%                           outside the axes (default none)
%   fill                    fill / shading colour
%   color                   colour of vectors, points, intervals
%   opacity                 fill opacity of planes and parallelepipeds (0.5)
%   style                   extra TikZ options for the main lines, e.g. dashed
%   scale                   figure scale (ellipse/ellipsoid default 1.5, else 1)
%   Number line
%   min, max, step          range and tick spacing (default -3, 3, 1)
%   tick labels=true|false  numbers under the ticks
%   points                  filled dots, e.g. {-1/$a$, 2.5/$b$} (value/label;
%                           write 2.5/{} for no label)
%   open points             open circles, same format
%   interval={from}{to}     highlight a segment, e.g. interval={-1}{2}
%   2D axes
%   xmin, xmax, ymin, ymax  axis ranges (default -1, 3, -1, 3)
%   step                    tick / grid spacing (default 1)
%   grid=true|false         light grid (default false)
%   tick labels=true|false  numbers along the axes (default false)
%   3D axes
%   min, max                axis range (default 0, 2.5)
%   Vectors, parallelograms, planes, parallelepipeds
%   Basis grid
%   u, v                    grid directions (default e1, e2); the grid lines
%                           are o + i*u + t*v and o + s*u + j*v for integers
%                           i, j. For the image of the standard grid under A,
%                           use u = A e1, v = A e2.
%   xmin, xmax, ymin, ymax  index ranges along u and along v (default -2..2)
%   color, style            line colour (default gray!60), extra options
%   vectors=true|false      draw u and v as arrows (default false), labelled
%                           with u label, v label
%   Vectors, parallelograms, planes, parallelepipeds
%   from, to                vector tail and tip, e.g. to={2,1} or to={1,2,1}
%   label                   vector label (at the tip by default)
%   label pos=tip|left|right
%                           left/right put the label beside the middle of
%                           the arrow (as seen walking from tail to tip);
%                           use when the tip is crowded
%   u, v, w                 spanning vectors, e.g. u={2,1} or u={1,0,1}
%   u label, v label, w label, sum label
%                           labels at the tips of u, v, w and at u+v
%   vectors=true|false      draw u, v (, w) as arrows (plane default false)
%   extent                  plane patch runs from -extent to +extent along
%                           u and v (default 1)
%   name pos={s}{t}         where a plane's name goes: s*extent along u plus
%                           t*extent along v (default {-0.7}{-0.7}); values
%                           beyond +/-1 put it just outside the patch
%   Ellipses, ellipsoids
%   a, b, c                 semi-axis lengths (numbers, default 1)
%   a label, b label, c label
%                           labels at the axis intercepts (default none)
%   y extent                make the figure reach at least this far above and
%                           below the origin; in 2D it is also the length of the
%                           x2 half-axis (default: just fit the figure). Give
%                           side-by-side figures the same y extent so their
%                           axes and above/below labels line up.
%   overhang                axis length beyond the intercepts (default 0.6)
%   azimuth, elevation      view for ellipsoids (default 30, 20; inside an
%                           llfig view picture, that picture's view)
%   equator=true|false      draw the equator x3 = 0 on ellipsoids (default true)
% ------------------------------------------------------------
\usepackage{tikz-3dplot}
\usetikzlibrary{calc}
\makeatletter
\newif\ifllfig@equator
\newif\ifllfig@ticks
\newif\ifllfig@grid
\newif\ifllfig@vectors
\newif\ifllfig@dimcap
\newif\ifllfig@dims
\newcount\llfig@k
\pgfkeys{
  /llfig/.is family, /llfig,
  a/.store in=\llfig@a,              b/.store in=\llfig@b,
  c/.store in=\llfig@c,
  x1 label/.store in=\llfig@xilabel, x2 label/.store in=\llfig@xiilabel,
  x3 label/.store in=\llfig@xiiilabel,
  a label/.store in=\llfig@alabel,   b label/.store in=\llfig@blabel,
  c label/.store in=\llfig@clabel,
  name/.store in=\llfig@name,
  label above/.store in=\llfig@above, label below/.store in=\llfig@below,
  y extent/.store in=\llfig@yext,
  fill/.store in=\llfig@fill,
  color/.store in=\llfig@color,
  opacity/.store in=\llfig@opacity,
  style/.code={\tikzset{llfig@userstyle/.style={#1}}},
  overhang/.store in=\llfig@over,
  scale/.store in=\llfig@scale,
  azimuth/.store in=\llfig@az,       elevation/.store in=\llfig@el,
  equator/.is if=llfig@equator,
  min/.store in=\llfig@min,          max/.store in=\llfig@max,
  step/.store in=\llfig@step,
  tick labels/.is if=llfig@ticks,
  points/.store in=\llfig@points,    open points/.store in=\llfig@opoints,
  interval/.code 2 args={\def\llfig@ivfrom{#1}\def\llfig@ivto{#2}},
  xmin/.store in=\llfig@xmin,        xmax/.store in=\llfig@xmax,
  ymin/.store in=\llfig@ymin,        ymax/.store in=\llfig@ymax,
  grid/.is if=llfig@grid,
  from/.store in=\llfig@from,        to/.store in=\llfig@to,
  label/.store in=\llfig@label,
  label pos/.store in=\llfig@lpos,
  before label/.store in=\llfig@beforelabel,
  after label/.store in=\llfig@afterlabel,
  u/.store in=\llfig@u,  v/.store in=\llfig@v,  w/.store in=\llfig@w,
  u label/.store in=\llfig@ulabel,   v label/.store in=\llfig@vlabel,
  w label/.store in=\llfig@wlabel,   sum label/.store in=\llfig@sumlabel,
  vectors/.is if=llfig@vectors,
  extent/.store in=\llfig@extent,
  window/.store in=\llfig@window,
  ambient/.store in=\llfig@ambient,
  vector name/.store in=\llfig@vname,
  dim caption/.is if=llfig@dimcap,
  domain/.store in=\llfig@domain,     codomain/.store in=\llfig@codomain,
  kernel/.store in=\llfig@kernel,     range/.store in=\llfig@range,
  map label/.store in=\llfig@maplabel,
  point/.store in=\llfig@ptlabel,     image/.store in=\llfig@imlabel,
  row space/.store in=\llfig@rowsp,   left null/.store in=\llfig@leftnul,
  dims/.is if=llfig@dims,
  rank/.store in=\llfig@rk,
  matrix name/.store in=\llfig@mname,
  name pos/.code 2 args={\def\llfig@ns{#1}\def\llfig@nt{#2}},
  reset/.style={a=1, b=1, c=1,
    x1 label=$x_1$, x2 label=$x_2$, x3 label=$x_3$,
    a label=, b label=, c label=, name=,
    label above=, label below=, y extent=,
    fill=gray!25, color=black, opacity=0.5, style={},
    overhang=0.6, scale=1.5,
    azimuth=30, elevation=20, equator=true,
    min=-3, max=3, step=1, tick labels=false,
    points=, open points=, interval={}{},
    xmin=-1, xmax=3, ymin=-1, ymax=3, grid=false,
    from={0,0}, to={1,0}, label=, label pos=tip,
    u={1,0}, v={0,1}, w={0,0,1},
    u label=, v label=, w label=, sum label=,
    vectors=true, extent=1, name pos={-0.7}{-0.7}},
}
% \llfig@ifnonempty{<macro>}{<code>}: run <code> only if <macro> is nonempty
% (expandable, so it also works in the middle of a TikZ path)
\newcommand{\llfig@ifnonempty}[2]{%
  \if\relax\detokenize\expandafter{#1}\relax\else#2\fi}
% \llfig@captions[<bounding box>]: draws label above / label below around
% a local bounding box (default llfig@bb, the one each pic sets up),
% centred horizontally on the point \llfig@captionx (default the origin).
% The fixed text height and depth keep labels of neighbouring figures on
% a common baseline (e.g. $D$ next to $E$).
\def\llfig@captionx{0,0}
\newcommand{\llfig@captions}[1][llfig@bb]{%
  \coordinate (llfig@origin) at (\llfig@captionx);%
  \llfig@ifnonempty\llfig@above{%
    \node[anchor=south, yshift=4pt, inner sep=0pt,
          text height=1.6ex, text depth=0.4ex]
      at (#1.north -| llfig@origin) {\llfig@above};}%
  \llfig@ifnonempty\llfig@below{%
    \node[anchor=north, yshift=-4pt, inner sep=0pt,
          text height=1.6ex, text depth=0.4ex]
      at (#1.south -| llfig@origin) {\llfig@below};}}
% \llfig@tiplabel{<tail>}{<tip>}{<node options>}{<macro>}: if <macro> is
% nonempty, put it just beyond the tip, on the far side from the tail.
\newcommand{\llfig@tiplabel}[4]{%
  \llfig@ifnonempty#4{%
    \path let \p1 = ($(#2)-(#1)$), \n1 = {atan2(\y1,\x1)+180} in
      node[anchor=\n1, inner sep=2pt, #3] at (#2) {#4};}}
% \llfig@splitxyz{<value>}{<name>}: split "x,y,z" into \llfig@<name>x, ...y, ...z
\def\llfig@splitxyz#1#2{\expandafter\llfig@splitxyz@i#1\relax{#2}}
\def\llfig@splitxyz@i#1,#2,#3\relax#4{%
  \expandafter\def\csname llfig@#4x\endcsname{#1}%
  \expandafter\def\csname llfig@#4y\endcsname{#2}%
  \expandafter\def\csname llfig@#4z\endcsname{#3}}
% \llfig@splitxy{<value>}{<name>}: split "x,y" into \llfig@<name>x, ...y
\def\llfig@splitxy#1#2{\expandafter\llfig@splitxy@i#1\relax{#2}}
\def\llfig@splitxy@i#1,#2\relax#3{%
  \expandafter\def\csname llfig@#3x\endcsname{#1}%
  \expandafter\def\csname llfig@#3y\endcsname{#2}}
% Linear maps, by changing the unit vectors: afterwards the coordinate x
% lands where A x used to. Both (all three) new units are computed from
% the old ones before any is set, so nested maps compose correctly.
% The combined 3x3 matrix in effect, M (identity outside any llfig
% matrix scope), is tracked so that pics needing true 3D directions (the
% parallelepiped's hidden-edge test) can use M u instead of u.
\def\llfig@Maa{1}\def\llfig@Mab{0}\def\llfig@Mac{0}
\def\llfig@Mba{0}\def\llfig@Mbb{1}\def\llfig@Mbc{0}
\def\llfig@Mca{0}\def\llfig@Mcb{0}\def\llfig@Mcc{1}
% \llfig@composeM: M := M * A, A's entries in \llfig@r<row><col> (rows a,
% b, c; columns x, y, z). Plain local definitions, so M stays inside
% the scope that set it.
\def\llfig@composeM{%
  \pgfmathsetmacro\llfig@Naa{\llfig@Maa*\llfig@rax + \llfig@Mab*\llfig@rbx + \llfig@Mac*\llfig@rcx}%
  \pgfmathsetmacro\llfig@Nab{\llfig@Maa*\llfig@ray + \llfig@Mab*\llfig@rby + \llfig@Mac*\llfig@rcy}%
  \pgfmathsetmacro\llfig@Nac{\llfig@Maa*\llfig@raz + \llfig@Mab*\llfig@rbz + \llfig@Mac*\llfig@rcz}%
  \pgfmathsetmacro\llfig@Nba{\llfig@Mba*\llfig@rax + \llfig@Mbb*\llfig@rbx + \llfig@Mbc*\llfig@rcx}%
  \pgfmathsetmacro\llfig@Nbb{\llfig@Mba*\llfig@ray + \llfig@Mbb*\llfig@rby + \llfig@Mbc*\llfig@rcy}%
  \pgfmathsetmacro\llfig@Nbc{\llfig@Mba*\llfig@raz + \llfig@Mbb*\llfig@rbz + \llfig@Mbc*\llfig@rcz}%
  \pgfmathsetmacro\llfig@Nca{\llfig@Mca*\llfig@rax + \llfig@Mcb*\llfig@rbx + \llfig@Mcc*\llfig@rcx}%
  \pgfmathsetmacro\llfig@Ncb{\llfig@Mca*\llfig@ray + \llfig@Mcb*\llfig@rby + \llfig@Mcc*\llfig@rcy}%
  \pgfmathsetmacro\llfig@Ncc{\llfig@Mca*\llfig@raz + \llfig@Mcb*\llfig@rbz + \llfig@Mcc*\llfig@rcz}%
  \let\llfig@Maa\llfig@Naa\let\llfig@Mab\llfig@Nab\let\llfig@Mac\llfig@Nac%
  \let\llfig@Mba\llfig@Nba\let\llfig@Mbb\llfig@Nbb\let\llfig@Mbc\llfig@Nbc%
  \let\llfig@Mca\llfig@Nca\let\llfig@Mcb\llfig@Ncb\let\llfig@Mcc\llfig@Ncc}
% \llfig@setmatrix{a,b}{c,d}: A = [a b; c d], given by rows
\def\llfig@setmatrix#1#2{%
  \llfig@splitxy{#1}{ra}\llfig@splitxy{#2}{rb}%
  \def\llfig@raz{0}\def\llfig@rbz{0}%
  \def\llfig@rcx{0}\def\llfig@rcy{0}\def\llfig@rcz{1}%
  \llfig@composeM
  \pgfpointxy{\llfig@rax}{\llfig@rbx}% A e1 = (a, c)
  \edef\llfig@ex{\the\pgf@x}\edef\llfig@ey{\the\pgf@y}%
  \pgfpointxy{\llfig@ray}{\llfig@rby}% A e2 = (b, d)
  \edef\llfig@fx{\the\pgf@x}\edef\llfig@fy{\the\pgf@y}%
  \pgfsetxvec{\pgfpoint{\llfig@ex}{\llfig@ey}}%
  \pgfsetyvec{\pgfpoint{\llfig@fx}{\llfig@fy}}}
% \llfig@setmatrixiii{row 1}{row 2}{row 3}: 3x3 A, rows like {a,b,c}
\def\llfig@setmatrixiii#1#2#3{%
  \llfig@splitxyz{#1}{ra}\llfig@splitxyz{#2}{rb}\llfig@splitxyz{#3}{rc}%
  \llfig@composeM
  \pgfpointxyz{\llfig@rax}{\llfig@rbx}{\llfig@rcx}% A e1
  \edef\llfig@ex{\the\pgf@x}\edef\llfig@ey{\the\pgf@y}%
  \pgfpointxyz{\llfig@ray}{\llfig@rby}{\llfig@rcy}% A e2
  \edef\llfig@fx{\the\pgf@x}\edef\llfig@fy{\the\pgf@y}%
  \pgfpointxyz{\llfig@raz}{\llfig@rbz}{\llfig@rcz}% A e3
  \edef\llfig@gx{\the\pgf@x}\edef\llfig@gy{\the\pgf@y}%
  \pgfsetxvec{\pgfpoint{\llfig@ex}{\llfig@ey}}%
  \pgfsetyvec{\pgfpoint{\llfig@fx}{\llfig@fy}}%
  \pgfsetzvec{\pgfpoint{\llfig@gx}{\llfig@gy}}}
% Inside an llfig view picture, ellipsoids default to that view at scale 1.
\newcommand{\llfig@useview}{%
  \ifdefined\llfig@viewaz
    \pgfkeys{/llfig/azimuth=\llfig@viewaz, /llfig/elevation=\llfig@viewel,
             /llfig/scale=1}%
  \fi}

\tikzset{
  llfig intercept/.style={inner sep=1.5pt},
  llfig@userstyle/.style={},
  % 3D view for a whole tikzpicture: llfig view={azimuth}{elevation}
  llfig view/.code 2 args={%
    \pgfmathsetmacro\llfig@vtheta{90 - (#2)}%
    \pgfmathsetmacro\llfig@vphi{90 + (#1)}%
    \tdplotsetmaincoords{\llfig@vtheta}{\llfig@vphi}%
    \def\llfig@viewaz{#1}\def\llfig@viewel{#2}%
    \tikzset{tdplot_main_coords}},
  llfig view/.default={30}{20},
  % scale the coordinate units (2D or 3D); pics follow, text does not
  llfig scale/.style={x={(#1,0,0)}, y={(0,#1,0)}, z={(0,0,#1)}},
  % apply x -> A x to everything drawn in the scope (rows of A)
  llfig matrix/.code 2 args={\llfig@setmatrix{#1}{#2}},
  llfig matrix 3d/.code n args={3}{\llfig@setmatrixiii{#1}{#2}{#3}},
  % ---------------- 2D: number line ----------------
  pics/number line/.style={code={%
    \pgfkeys{/llfig, reset, scale=1, x1 label=, color=blue,
             tick labels=true, #1}%
    % scale the unit only, so dots and tick marks keep their size
    \begin{scope}[x={(\llfig@scale,0)}, >=stealth]
      \begin{scope}[local bounding box=llfig@bb]
        \draw[<->, llfig@userstyle] ({\llfig@min-0.5},0) -- ({\llfig@max+0.5},0)
          \llfig@ifnonempty\llfig@xilabel{node[right] {\llfig@xilabel}};
        \pgfmathsetmacro\llfig@next{\llfig@min + \llfig@step}
        \foreach \x in {\llfig@min,\llfig@next,...,\llfig@max} {
          \draw (\x,3pt) -- (\x,-3pt);
          \ifllfig@ticks
            \node[below=4pt, inner sep=1pt] at (\x,0) {$\pgfmathprintnumber{\x}$};
          \fi
        }
        \llfig@ifnonempty\llfig@ivfrom{%
          \draw[line width=3pt, color=\llfig@color, opacity=0.5]
            (\llfig@ivfrom,0) -- (\llfig@ivto,0);}
        \llfig@ifnonempty\llfig@points{%
          \foreach \p/\l in \llfig@points {
            \fill[color=\llfig@color] (\p,0) circle (2.5pt)
              node[above=3pt, text=black] {\l};}}
        \llfig@ifnonempty\llfig@opoints{%
          \foreach \p/\l in \llfig@opoints {
            \filldraw[draw=\llfig@color, fill=white, thick] (\p,0) circle (2.5pt)
              node[above=3pt, text=black] {\l};}}
      \end{scope}
      \def\llfig@captionx{llfig@bb.center}%
      \llfig@captions
    \end{scope}
  }},
  pics/number line/.default=,
  % ---------------- 2D: axes ----------------
  pics/axes 2d/.style={code={%
    \pgfkeys{/llfig, reset, scale=1, #1}%
    \begin{scope}[x={(\llfig@scale,0)}, y={(0,\llfig@scale)}, >=stealth]
      \begin{scope}[local bounding box=llfig@bb]
        \ifllfig@grid
          \draw[gray!30, very thin, step=\llfig@step]
            (\llfig@xmin,\llfig@ymin) grid (\llfig@xmax,\llfig@ymax);
        \fi
        \draw[->] ({\llfig@xmin-0.4},0) -- ({\llfig@xmax+0.4},0)
          \llfig@ifnonempty\llfig@xilabel{node[right] {\llfig@xilabel}};
        \draw[->] (0,{\llfig@ymin-0.4}) -- (0,{\llfig@ymax+0.4})
          \llfig@ifnonempty\llfig@xiilabel{node[above] {\llfig@xiilabel}};
        \ifllfig@ticks
          \pgfmathsetmacro\llfig@next{\llfig@xmin + \llfig@step}
          \foreach \x in {\llfig@xmin,\llfig@next,...,\llfig@xmax} {
            \ifdim \x pt=0pt \else
              \draw (\x,2pt) -- (\x,-2pt)
                node[below, inner sep=2pt] {\scriptsize$\pgfmathprintnumber{\x}$};
            \fi}
          \pgfmathsetmacro\llfig@next{\llfig@ymin + \llfig@step}
          \foreach \y in {\llfig@ymin,\llfig@next,...,\llfig@ymax} {
            \ifdim \y pt=0pt \else
              \draw (2pt,\y) -- (-2pt,\y)
                node[left, inner sep=2pt] {\scriptsize$\pgfmathprintnumber{\y}$};
            \fi}
        \fi
      \end{scope}
      \llfig@captions
    \end{scope}
  }},
  pics/axes 2d/.default=,
  % ---------------- 2D or 3D: vector ----------------
  pics/vector/.style={code={%
    \pgfkeys{/llfig, reset, #1}%
    \coordinate (llfig@from) at (\llfig@from);
    \coordinate (llfig@to) at (\llfig@to);
    \def\llfig@tmp{tip}%
    \edef\llfig@lpos{\llfig@lpos}% (the value may arrive as a macro)
    \ifx\llfig@lpos\llfig@tmp
      \draw[->, >=stealth, thick, color=\llfig@color, llfig@userstyle]
        (llfig@from) -- (llfig@to);
      \llfig@tiplabel{llfig@from}{llfig@to}{text=\llfig@color}{\llfig@label}
    \else % label pos=left|right: beside the middle of the arrow
      \draw[->, >=stealth, thick, color=\llfig@color, llfig@userstyle]
        (llfig@from) -- (llfig@to)
        node[pos=0.5, auto=\llfig@lpos, inner sep=2pt] {\llfig@label};
    \fi
  }},
  pics/vector/.default=,
  % ---------------- 2D or 3D: parallelogram ----------------
  pics/parallelogram/.style={code={%
    \pgfkeys{/llfig, reset, u={2,1}, v={1,2}, fill=blue!20, opacity=1, #1}%
    \coordinate (llfig@o) at (0,0);
    \coordinate (llfig@u) at (\llfig@u);
    \coordinate (llfig@v) at (\llfig@v);
    \coordinate (llfig@w) at ($(llfig@u) + (llfig@v) - (llfig@o)$);
    \begin{scope}[local bounding box=llfig@bb, >=stealth]
      \filldraw[fill=\llfig@fill, fill opacity=\llfig@opacity, draw=black,
                llfig@userstyle]
        (llfig@o) -- (llfig@u) -- (llfig@w) -- (llfig@v) -- cycle;
      \ifllfig@vectors
        \draw[->, thick, color=\llfig@color] (llfig@o) -- (llfig@u);
        \draw[->, thick, color=\llfig@color] (llfig@o) -- (llfig@v);
      \fi
      % corner labels point away from the centre, clear of both edges
      \coordinate (llfig@c) at ($(llfig@o)!0.5!(llfig@w)$);
      \llfig@tiplabel{llfig@c}{llfig@u}{}{\llfig@ulabel}
      \llfig@tiplabel{llfig@c}{llfig@v}{}{\llfig@vlabel}
      \llfig@tiplabel{llfig@c}{llfig@w}{}{\llfig@sumlabel}
      \llfig@ifnonempty\llfig@name{\node at (llfig@c) {\llfig@name};}
    \end{scope}
    \llfig@captions
  }},
  pics/parallelogram/.default=,
  % ---------------- 2D or 3D: grid along a basis ----------------
  % Lines o + i*(u-o) + t*(v-o), t in [ymin, ymax], for integer i in
  % [xmin, xmax], and likewise along v. (Coordinates already include the
  % pic's shift o, hence the multiples of (llfig@o).)
  pics/basis grid/.style={code={%
    \pgfkeys{/llfig, reset, color=gray!60, vectors=false,
             xmin=-2, xmax=2, ymin=-2, ymax=2, #1}%
    \coordinate (llfig@o) at (0,0);
    \coordinate (llfig@u) at (\llfig@u);
    \coordinate (llfig@v) at (\llfig@v);
    \begin{scope}[local bounding box=llfig@bb, >=stealth]
      \foreach \i in {\llfig@xmin,...,\llfig@xmax} {
        \draw[color=\llfig@color, llfig@userstyle]
          ($(llfig@o) + {\i}*(llfig@u) + {\llfig@ymin}*(llfig@v)
             - {\i + \llfig@ymin}*(llfig@o)$)
          -- ($(llfig@o) + {\i}*(llfig@u) + {\llfig@ymax}*(llfig@v)
             - {\i + \llfig@ymax}*(llfig@o)$);}
      \foreach \j in {\llfig@ymin,...,\llfig@ymax} {
        \draw[color=\llfig@color, llfig@userstyle]
          ($(llfig@o) + {\llfig@xmin}*(llfig@u) + {\j}*(llfig@v)
             - {\llfig@xmin + \j}*(llfig@o)$)
          -- ($(llfig@o) + {\llfig@xmax}*(llfig@u) + {\j}*(llfig@v)
             - {\llfig@xmax + \j}*(llfig@o)$);}
      \ifllfig@vectors
        \draw[->, thick] (llfig@o) -- (llfig@u);
        \draw[->, thick] (llfig@o) -- (llfig@v);
        % labels beside the middle of each arrow, on the side away from the
        % other vector, so they sit inside a grid cell rather than on a line
        \llfig@ifnonempty\llfig@ulabel{%
          \path let \p1 = ($(llfig@v)-(llfig@o)$), \n1 = {atan2(\y1,\x1)} in
            node[anchor=\n1, inner sep=2pt] at ($(llfig@o)!0.5!(llfig@u)$)
            {\llfig@ulabel};}
        \llfig@ifnonempty\llfig@vlabel{%
          \path let \p1 = ($(llfig@u)-(llfig@o)$), \n1 = {atan2(\y1,\x1)} in
            node[anchor=\n1, inner sep=2pt] at ($(llfig@o)!0.5!(llfig@v)$)
            {\llfig@vlabel};}
      \fi
    \end{scope}
    \llfig@captions
  }},
  pics/basis grid/.default=,
  % ---------------- 3D: axes ----------------
  pics/axes 3d/.style={code={%
    \pgfkeys{/llfig, reset, min=0, max=2.5, #1}%
    \coordinate (llfig@o) at (0,0,0);
    \coordinate (llfig@e1) at (\llfig@max,0,0);
    \coordinate (llfig@e2) at (0,\llfig@max,0);
    \coordinate (llfig@e3) at (0,0,\llfig@max);
    \begin{scope}[local bounding box=llfig@bb, >=stealth]
      \draw[->] (\llfig@min,0,0) -- (llfig@e1);
      \draw[->] (0,\llfig@min,0) -- (llfig@e2);
      \draw[->] (0,0,\llfig@min) -- (llfig@e3);
      \llfig@tiplabel{llfig@o}{llfig@e1}{}{\llfig@xilabel}
      \llfig@tiplabel{llfig@o}{llfig@e2}{}{\llfig@xiilabel}
      \llfig@tiplabel{llfig@o}{llfig@e3}{}{\llfig@xiiilabel}
    \end{scope}
    \llfig@captions
  }},
  pics/axes 3d/.default=,
  % ---------------- 3D: plane patch ----------------
  % Corners o +/- extent*u +/- extent*v. (Coordinates already include the
  % pic's shift o, hence the multiples of (llfig@o) below.)
  pics/plane/.style={code={%
    \pgfkeys{/llfig, reset, u={1,0,0}, v={0,1,0}, fill=blue!25,
             vectors=false, #1}%
    \coordinate (llfig@o) at (0,0,0);
    \coordinate (llfig@u) at (\llfig@u);
    \coordinate (llfig@v) at (\llfig@v);
    \begin{scope}[local bounding box=llfig@bb, >=stealth]
      \filldraw[fill=\llfig@fill, fill opacity=\llfig@opacity, draw=black!60,
                llfig@userstyle]
        ($(llfig@o) - {\llfig@extent}*(llfig@u) - {\llfig@extent}*(llfig@v)
           + {2*\llfig@extent}*(llfig@o)$)
        -- ($(llfig@o) + {\llfig@extent}*(llfig@u) - {\llfig@extent}*(llfig@v)$)
        -- ($(llfig@o) + {\llfig@extent}*(llfig@u) + {\llfig@extent}*(llfig@v)
              - {2*\llfig@extent}*(llfig@o)$)
        -- ($(llfig@o) - {\llfig@extent}*(llfig@u) + {\llfig@extent}*(llfig@v)$)
        -- cycle;
      \ifllfig@vectors
        \draw[->, thick, color=\llfig@color] (llfig@o) -- (llfig@u);
        \draw[->, thick, color=\llfig@color] (llfig@o) -- (llfig@v);
        \llfig@tiplabel{llfig@o}{llfig@u}{}{\llfig@ulabel}
        \llfig@tiplabel{llfig@o}{llfig@v}{}{\llfig@vlabel}
      \fi
      % name at o + s*extent*(u - o) + t*extent*(v - o), (s,t) = name pos
      \llfig@ifnonempty\llfig@name{%
        \node at ($(llfig@o) + {\llfig@ns*\llfig@extent}*(llfig@u)
          + {\llfig@nt*\llfig@extent}*(llfig@v)
          - {(\llfig@ns + \llfig@nt)*\llfig@extent}*(llfig@o)$)
          {\llfig@name};}
    \end{scope}
    \llfig@captions
  }},
  pics/plane/.default=,
  % ---------------- 3D: parallelepiped ----------------
  % Corner p<i><j><k> = i*u + j*v + k*w. Writing the direction toward the
  % viewer as d = alpha*u + beta*v + gamma*w (Cramer's rule), the corner
  % farthest from the viewer has i = 1 iff alpha < 0, etc.; exactly the
  % three edges at that corner are hidden, and the three faces not
  % touching it face the viewer.
  pics/parallelepiped/.style={code={%
    \pgfkeys{/llfig, reset, u={2,0,0}, v={0,2,0}, w={0,0,2},
             fill=blue!25, #1}%
    \ifdefined\rcarot\else\tdplotsetmaincoords{70}{120}\fi
    \llfig@splitxyz\llfig@u{u}\llfig@splitxyz\llfig@v{v}%
    \llfig@splitxyz\llfig@w{w}%
    % The test needs the true directions: M u, M v, M w, where M is the
    % combined llfig matrix in effect (identity if none). Drawing still
    % uses u, v, w, which the matrix then maps.
    \pgfmathsetmacro\llfig@Wux{\llfig@Maa*\llfig@ux + \llfig@Mab*\llfig@uy + \llfig@Mac*\llfig@uz}%
    \pgfmathsetmacro\llfig@Wuy{\llfig@Mba*\llfig@ux + \llfig@Mbb*\llfig@uy + \llfig@Mbc*\llfig@uz}%
    \pgfmathsetmacro\llfig@Wuz{\llfig@Mca*\llfig@ux + \llfig@Mcb*\llfig@uy + \llfig@Mcc*\llfig@uz}%
    \pgfmathsetmacro\llfig@Wvx{\llfig@Maa*\llfig@vx + \llfig@Mab*\llfig@vy + \llfig@Mac*\llfig@vz}%
    \pgfmathsetmacro\llfig@Wvy{\llfig@Mba*\llfig@vx + \llfig@Mbb*\llfig@vy + \llfig@Mbc*\llfig@vz}%
    \pgfmathsetmacro\llfig@Wvz{\llfig@Mca*\llfig@vx + \llfig@Mcb*\llfig@vy + \llfig@Mcc*\llfig@vz}%
    \pgfmathsetmacro\llfig@Wwx{\llfig@Maa*\llfig@wx + \llfig@Mab*\llfig@wy + \llfig@Mac*\llfig@wz}%
    \pgfmathsetmacro\llfig@Wwy{\llfig@Mba*\llfig@wx + \llfig@Mbb*\llfig@wy + \llfig@Mbc*\llfig@wz}%
    \pgfmathsetmacro\llfig@Wwz{\llfig@Mca*\llfig@wx + \llfig@Mcb*\llfig@wy + \llfig@Mcc*\llfig@wz}%
    % sign tests: (d . (v x w)) * det[u v w] < 0  <=>  alpha < 0, etc.
    \pgfmathsetmacro\llfig@det{\llfig@Wux*(\llfig@Wvy*\llfig@Wwz - \llfig@Wvz*\llfig@Wwy)
      + \llfig@Wuy*(\llfig@Wvz*\llfig@Wwx - \llfig@Wvx*\llfig@Wwz)
      + \llfig@Wuz*(\llfig@Wvx*\llfig@Wwy - \llfig@Wvy*\llfig@Wwx)}%
    \pgfmathtruncatemacro\llfig@bu{\llfig@det*(
        \rcarot*(\llfig@Wvy*\llfig@Wwz - \llfig@Wvz*\llfig@Wwy)
      + \rcbrot*(\llfig@Wvz*\llfig@Wwx - \llfig@Wvx*\llfig@Wwz)
      + \rccrot*(\llfig@Wvx*\llfig@Wwy - \llfig@Wvy*\llfig@Wwx)) < 0}%
    \pgfmathtruncatemacro\llfig@bv{\llfig@det*(
        \rcarot*(\llfig@Wwy*\llfig@Wuz - \llfig@Wwz*\llfig@Wuy)
      + \rcbrot*(\llfig@Wwz*\llfig@Wux - \llfig@Wwx*\llfig@Wuz)
      + \rccrot*(\llfig@Wwx*\llfig@Wuy - \llfig@Wwy*\llfig@Wux)) < 0}%
    \pgfmathtruncatemacro\llfig@bw{\llfig@det*(
        \rcarot*(\llfig@Wuy*\llfig@Wvz - \llfig@Wuz*\llfig@Wvy)
      + \rcbrot*(\llfig@Wuz*\llfig@Wvx - \llfig@Wux*\llfig@Wvz)
      + \rccrot*(\llfig@Wux*\llfig@Wvy - \llfig@Wuy*\llfig@Wvx)) < 0}%
    \pgfmathtruncatemacro\llfig@fu{1 - \llfig@bu}%
    \pgfmathtruncatemacro\llfig@fv{1 - \llfig@bv}%
    \pgfmathtruncatemacro\llfig@fw{1 - \llfig@bw}%
    \foreach \i in {0,1} {\foreach \j in {0,1} {\foreach \k in {0,1} {
      \coordinate (llfig@p\i\j\k) at
        ({\i*\llfig@ux + \j*\llfig@vx + \k*\llfig@wx},
         {\i*\llfig@uy + \j*\llfig@vy + \k*\llfig@wy},
         {\i*\llfig@uz + \j*\llfig@vz + \k*\llfig@wz});}}}
    \begin{scope}[local bounding box=llfig@bb, >=stealth]
      % front faces
      \fill[fill=\llfig@fill, fill opacity=\llfig@opacity]
        (llfig@p\llfig@fu00) -- (llfig@p\llfig@fu10)
        -- (llfig@p\llfig@fu11) -- (llfig@p\llfig@fu01) -- cycle;
      \fill[fill=\llfig@fill, fill opacity=\llfig@opacity]
        (llfig@p0\llfig@fv0) -- (llfig@p1\llfig@fv0)
        -- (llfig@p1\llfig@fv1) -- (llfig@p0\llfig@fv1) -- cycle;
      \fill[fill=\llfig@fill, fill opacity=\llfig@opacity]
        (llfig@p00\llfig@fw) -- (llfig@p10\llfig@fw)
        -- (llfig@p11\llfig@fw) -- (llfig@p01\llfig@fw) -- cycle;
      % edges: dashed iff they meet the back corner
      \foreach \j in {0,1} {\foreach \k in {0,1} {
        \pgfmathtruncatemacro\llfig@h{and(\j == \llfig@bv, \k == \llfig@bw)}
        \ifnum\llfig@h=1 \tikzset{llfig@edge/.style={dashed}}%
        \else \tikzset{llfig@edge/.style={}}\fi
        \draw[llfig@edge, llfig@userstyle] (llfig@p0\j\k) -- (llfig@p1\j\k);
        \pgfmathtruncatemacro\llfig@h{and(\j == \llfig@bu, \k == \llfig@bw)}
        \ifnum\llfig@h=1 \tikzset{llfig@edge/.style={dashed}}%
        \else \tikzset{llfig@edge/.style={}}\fi
        \draw[llfig@edge, llfig@userstyle] (llfig@p\j0\k) -- (llfig@p\j1\k);
        \pgfmathtruncatemacro\llfig@h{and(\j == \llfig@bu, \k == \llfig@bv)}
        \ifnum\llfig@h=1 \tikzset{llfig@edge/.style={dashed}}%
        \else \tikzset{llfig@edge/.style={}}\fi
        \draw[llfig@edge, llfig@userstyle] (llfig@p\j\k0) -- (llfig@p\j\k1);
      }}
      \ifllfig@vectors
        % an edge from the origin is hidden iff the back corner is on it
        \pgfmathtruncatemacro\llfig@h{and(\llfig@bv == 0, \llfig@bw == 0)}
        \ifnum\llfig@h=1 \tikzset{llfig@edge/.style={dashed}}%
        \else \tikzset{llfig@edge/.style={}}\fi
        \draw[->, thick, color=\llfig@color, llfig@edge]
          (llfig@p000) -- (llfig@p100);
        \pgfmathtruncatemacro\llfig@h{and(\llfig@bu == 0, \llfig@bw == 0)}
        \ifnum\llfig@h=1 \tikzset{llfig@edge/.style={dashed}}%
        \else \tikzset{llfig@edge/.style={}}\fi
        \draw[->, thick, color=\llfig@color, llfig@edge]
          (llfig@p000) -- (llfig@p010);
        \pgfmathtruncatemacro\llfig@h{and(\llfig@bu == 0, \llfig@bv == 0)}
        \ifnum\llfig@h=1 \tikzset{llfig@edge/.style={dashed}}%
        \else \tikzset{llfig@edge/.style={}}\fi
        \draw[->, thick, color=\llfig@color, llfig@edge]
          (llfig@p000) -- (llfig@p001);
      \fi
      % corner labels point away from the centre, so corners on the
      % outline get labels outside the solid
      \coordinate (llfig@c) at ($(llfig@p000)!0.5!(llfig@p111)$);
      \llfig@tiplabel{llfig@c}{llfig@p100}{}{\llfig@ulabel}
      \llfig@tiplabel{llfig@c}{llfig@p010}{}{\llfig@vlabel}
      \llfig@tiplabel{llfig@c}{llfig@p001}{}{\llfig@wlabel}
      \llfig@ifnonempty\llfig@name{\node at (llfig@c) {\llfig@name};}
    \end{scope}
    \llfig@captions
  }},
  pics/parallelepiped/.default=,
  % ---------------- 2D: ellipse / disk ----------------
  pics/ellipse/.style={code={%
    \pgfkeys{/llfig, reset, #1}%
    \begin{scope}[scale=\llfig@scale, >=stealth]
      \begin{scope}[local bounding box=llfig@bb]
        \pgfmathsetmacro\llfig@xmax{\llfig@a + \llfig@over}
        \pgfmathsetmacro\llfig@ymax{\llfig@b + \llfig@over}
        \llfig@ifnonempty\llfig@yext{\let\llfig@ymax\llfig@yext}
        \filldraw[fill=\llfig@fill, draw=black]
          (0,0) ellipse [x radius=\llfig@a, y radius=\llfig@b];
        \draw[->] (-\llfig@xmax,0) -- (\llfig@xmax,0)
          \llfig@ifnonempty\llfig@xilabel{node[right] {\llfig@xilabel}};
        \draw[->] (0,-\llfig@ymax) -- (0,\llfig@ymax)
          \llfig@ifnonempty\llfig@xiilabel{node[above] {\llfig@xiilabel}};
        \llfig@ifnonempty\llfig@alabel{\node[below right] at (\llfig@a,0) {\llfig@alabel};}
        \llfig@ifnonempty\llfig@blabel{\node[above left] at (0,\llfig@b) {\llfig@blabel};}
        \llfig@ifnonempty\llfig@name{%
          \node[above right] at ({\llfig@a*cos(45)},{\llfig@b*sin(45)}) {\llfig@name};}
      \end{scope}
      \llfig@captions
    \end{scope}
  }},
  pics/ellipse/.default=,
  % ---------------- 3D: ellipsoid / sphere ----------------
  pics/ellipsoid/.style={code={%
    \pgfkeys{/llfig, reset}\llfig@useview\pgfkeys{/llfig, #1}%
    \pgfmathsetmacro\llfig@theta{90 - \llfig@el}%
    \pgfmathsetmacro\llfig@phi{90 + \llfig@az}%
    \tdplotsetmaincoords{\llfig@theta}{\llfig@phi}%
    % tikz-3dplot's view rotation R: rows 1-2 (\raarot ... \rbcrot) map
    % (x1,x2,x3) to the page; row 3 (\rcarot ...) points at the viewer.
    % Outline = image of the unit circle under L, where L L^T = P D^2 P^T
    % (P = rows 1-2 of R, D = diag(a,b,c)); L is the 2x2 Cholesky factor.
    \pgfmathsetmacro\llfig@Sxx{(\llfig@a*\raarot)^2 + (\llfig@b*\rabrot)^2
                             + (\llfig@c*\racrot)^2}%
    \pgfmathsetmacro\llfig@Sxy{(\llfig@a)^2*\raarot*\rbarot + (\llfig@b)^2*\rabrot*\rbbrot
                             + (\llfig@c)^2*\racrot*\rbcrot}%
    \pgfmathsetmacro\llfig@Syy{(\llfig@a*\rbarot)^2 + (\llfig@b*\rbbrot)^2
                             + (\llfig@c*\rbcrot)^2}%
    \pgfmathsetmacro\llfig@Lxx{sqrt(\llfig@Sxx)}%
    \pgfmathsetmacro\llfig@Lyx{\llfig@Sxy/\llfig@Lxx}%
    \pgfmathsetmacro\llfig@Lyy{sqrt(\llfig@Syy - (\llfig@Lyx)^2)}%
    % The visible half of the equator is centred on angle alpha.
    \pgfmathsetmacro\llfig@alpha{atan2(\rcbrot/\llfig@b, \rcarot/\llfig@a)}%
    % Page length of a unit step along each axis; overhangs in world units
    % that all appear \llfig@over long on the page.
    \pgfmathsetmacro\llfig@ni{sqrt((\raarot)^2 + (\rbarot)^2)}%
    \pgfmathsetmacro\llfig@nii{sqrt((\rabrot)^2 + (\rbbrot)^2)}%
    \pgfmathsetmacro\llfig@niii{sqrt((\racrot)^2 + (\rbcrot)^2)}%
    \pgfmathsetmacro\llfig@oi{\llfig@over/\llfig@ni}%
    \pgfmathsetmacro\llfig@oii{\llfig@over/\llfig@nii}%
    \pgfmathsetmacro\llfig@oiii{\llfig@over/\llfig@niii}%
    % (the outline is drawn in page coordinates, even inside llfig view)
    \begin{scope}[x={(1cm,0cm)}, y={(0cm,1cm)}, scale=\llfig@scale, >=stealth]
      \begin{scope}[local bounding box=llfig@bb]
        \llfig@ifnonempty\llfig@yext{\path (0,-\llfig@yext) (0,\llfig@yext);}
        \shade[inner color=white, outer color=\llfig@fill, variable=\t,
               smooth, samples=73]
          plot[domain=0:360]
          ({\llfig@Lxx*cos(\t)}, {\llfig@Lyx*cos(\t) + \llfig@Lyy*sin(\t)}) -- cycle;
        \draw[variable=\t, smooth, samples=73]
          plot[domain=0:360]
          ({\llfig@Lxx*cos(\t)}, {\llfig@Lyx*cos(\t) + \llfig@Lyy*sin(\t)}) -- cycle;
        \begin{scope}[tdplot_main_coords]
          \ifllfig@equator
            \draw[variable=\t, smooth, samples=37]
              plot[domain={\llfig@alpha-90}:{\llfig@alpha+90}]
              ({\llfig@a*cos(\t)}, {\llfig@b*sin(\t)}, 0);
            \draw[dashed, variable=\t, smooth, samples=37]
              plot[domain={\llfig@alpha+90}:{\llfig@alpha+270}]
              ({\llfig@a*cos(\t)}, {\llfig@b*sin(\t)}, 0);
          \fi
          % hidden (inside) parts of the axes are dashed
          \draw[dashed] (0,0,0) -- (\llfig@a,0,0);
          \draw[dashed] (0,0,0) -- (0,\llfig@b,0);
          \draw[dashed] (0,0,0) -- (0,0,\llfig@c);
          \draw[->] (\llfig@a,0,0) -- ({\llfig@a+\llfig@oi},0,0)
            \llfig@ifnonempty\llfig@xilabel{node[below left] {\llfig@xilabel}};
          \draw[->] (0,\llfig@b,0) -- (0,{\llfig@b+\llfig@oii},0)
            \llfig@ifnonempty\llfig@xiilabel{node[right] {\llfig@xiilabel}};
          \draw[->] (0,0,\llfig@c) -- (0,0,{\llfig@c+\llfig@oiii})
            \llfig@ifnonempty\llfig@xiiilabel{node[above] {\llfig@xiiilabel}};
          % intercept labels sit 0.3cm (on the page) past the intercept,
          % beside the axis line, so they clear the outline and equator
          \llfig@ifnonempty\llfig@alabel{%
            \node[llfig intercept, anchor=north west]
              at ({\llfig@a + 0.3/\llfig@ni},0,0) {\llfig@alabel};}
          \llfig@ifnonempty\llfig@blabel{%
            \node[llfig intercept, anchor=north]
              at (0,{\llfig@b + 0.3/\llfig@nii},0) {\llfig@blabel};}
          \llfig@ifnonempty\llfig@clabel{%
            \node[llfig intercept, anchor=east]
              at (0,0,{\llfig@c + 0.3/\llfig@niii}) {\llfig@clabel};}
        \end{scope}
        \llfig@ifnonempty\llfig@name{%
          \node[above right] at ({\llfig@Lxx*cos(45)},
            {\llfig@Lyx*cos(45) + \llfig@Lyy*sin(45)}) {\llfig@name};}
      \end{scope}
      \llfig@captions
    \end{scope}
  }},
  pics/ellipsoid/.default=,
}

% Stand-alone shortcuts: one tikzpicture per figure, baseline at the origin.
\newcommand{\EllipseFig}[1][]{\tikz[baseline]\pic{ellipse={#1}};}
\newcommand{\UnitDiskFig}[1][]{\EllipseFig[a label=$1$, b label=$1$, #1]}
\newcommand{\EllipsoidFig}[1][]{\tikz[baseline]\pic{ellipsoid={#1}};}
\newcommand{\UnitSphereFig}[1][]{%
  \EllipsoidFig[a label=$1$, b label=$1$, c label=$1$, #1]}
\newcommand{\SpheroidFig}[1][]{%
  \begingroup\pgfkeys{/llfig, reset, #1}%
  \edef\llfig@tmp{\noexpand\EllipsoidFig[#1, b=\llfig@a]}%
  \expandafter\endgroup\llfig@tmp}
\newcommand{\NumberLineFig}[1][]{\tikz[baseline]\pic{number line={#1}};}
% Parallelogram on axes that just fit it. Captions go under the whole figure.
\newcommand{\ParallelogramFig}[1][]{%
  \begingroup
  \pgfkeys{/llfig, reset, scale=1, u={2,1}, v={1,2}, #1}%
  \llfig@splitxy\llfig@u{u}\llfig@splitxy\llfig@v{v}%
  \pgfmathsetmacro\llfig@pxmin{min(0, \llfig@ux, \llfig@vx, \llfig@ux + \llfig@vx) - 0.3}%
  \pgfmathsetmacro\llfig@pxmax{max(0, \llfig@ux, \llfig@vx, \llfig@ux + \llfig@vx) + 0.3}%
  \pgfmathsetmacro\llfig@pymin{min(0, \llfig@uy, \llfig@vy, \llfig@uy + \llfig@vy) - 0.3}%
  \pgfmathsetmacro\llfig@pymax{max(0, \llfig@uy, \llfig@vy, \llfig@uy + \llfig@vy) + 0.3}%
  \begin{tikzpicture}[baseline, llfig scale=\llfig@scale]
    \begin{scope}[local bounding box=llfig@whole]
      \pic {axes 2d={#1, scale=1, xmin=\llfig@pxmin, xmax=\llfig@pxmax,
                     ymin=\llfig@pymin, ymax=\llfig@pymax,
                     label above=, label below=}};
      \pic {parallelogram={#1, label above=, label below=}};
    \end{scope}
    \def\llfig@captionx{llfig@whole.center}%
    \llfig@captions[llfig@whole]
  \end{tikzpicture}%
  \endgroup}
% Parallelepiped on 3D axes that just fit it, viewed from azimuth/elevation.
\newcommand{\ParallelepipedFig}[1][]{%
  \begingroup
  \pgfkeys{/llfig, reset, scale=1, u={2,0,0}, v={0,2,0}, w={0,0,2}, #1}%
  \llfig@splitxyz\llfig@u{u}\llfig@splitxyz\llfig@v{v}\llfig@splitxyz\llfig@w{w}%
  \pgfmathsetmacro\llfig@pmax{max(\llfig@ux + \llfig@vx + \llfig@wx,
    \llfig@uy + \llfig@vy + \llfig@wy, \llfig@uz + \llfig@vz + \llfig@wz,
    \llfig@ux, \llfig@vy, \llfig@wz) + 1}%
  \begin{tikzpicture}[baseline, llfig view={\llfig@az}{\llfig@el},
                      llfig scale=\llfig@scale]
    \begin{scope}[local bounding box=llfig@whole]
      \pic {axes 3d={#1, max=\llfig@pmax, label above=, label below=}};
      \pic {parallelepiped={#1, label above=, label below=}};
    \end{scope}
    \def\llfig@captionx{llfig@whole.center}%
    \llfig@captions[llfig@whole]
  \end{tikzpicture}%
  \endgroup}
% Content, an arrow, and the same content under the matrix with rows
% #2 and #3. Both pictures have their baseline at the origin.
\newcommand{\MatrixMapFig}[4][]{%
  \begingroup
  \pgfkeys{/llfig, reset, scale=1, label=$T$, before label=, after label=, #1}%
  \begin{tikzpicture}[baseline, llfig scale=\llfig@scale]
    \begin{scope}[local bounding box=llfig@whole]#4\end{scope}
    \let\llfig@below\llfig@beforelabel
    \llfig@captions[llfig@whole]
  \end{tikzpicture}%
  \quad$\xrightarrow{\;\text{\llfig@label}\;}$\quad
  \begin{tikzpicture}[baseline, llfig scale=\llfig@scale,
                      llfig matrix={#2}{#3}]
    \begin{scope}[local bounding box=llfig@whole]#4\end{scope}
    \let\llfig@below\llfig@afterlabel
    \llfig@captions[llfig@whole]
  \end{tikzpicture}%
  \endgroup}

% ---------------- Span of columns in R^1, R^2 or R^3 ----------------
% Columns are stored as \llfig@V<i>x, ...y, ...z (missing entries 0).
\def\llfig@col#1#2{\csname llfig@V#1#2\endcsname}
\def\llfig@readvec#1,#2,#3,#4\relax{%
  \expandafter\gdef\csname llfig@V\the\llfig@k x\endcsname{#1}%
  \if\relax\detokenize{#2}\relax
    \expandafter\gdef\csname llfig@V\the\llfig@k y\endcsname{0}%
  \else\expandafter\gdef\csname llfig@V\the\llfig@k y\endcsname{#2}\fi
  \if\relax\detokenize{#3}\relax
    \expandafter\gdef\csname llfig@V\the\llfig@k z\endcsname{0}%
  \else\expandafter\gdef\csname llfig@V\the\llfig@k z\endcsname{#3}\fi
  \ifnum\llfig@k=1 % ambient dimension from the first column
    \if\relax\detokenize{#2}\relax \gdef\llfig@n{1}%
    \else\if\relax\detokenize{#3}\relax \gdef\llfig@n{2}%
    \else \gdef\llfig@n{3}\fi\fi
  \fi}
% Rank of the columns (tolerance 0.0001; pgfmath is fixed point), with witnesses: \llfig@p (a
% nonzero column; with \llfig@q, an independent pair).
\def\llfig@findrank{%
  \gdef\llfig@rank{0}\gdef\llfig@p{0}\gdef\llfig@q{0}%
  \ifnum\llfig@k>0
    \foreach \i in {1,...,\the\llfig@k} {%
      \pgfmathparse{(\llfig@col{\i}{x})^2 + (\llfig@col{\i}{y})^2
                    + (\llfig@col{\i}{z})^2 > 0.0001}%
      \ifnum\llfig@rank=0 \ifnum\pgfmathresult=1
        \gdef\llfig@rank{1}\xdef\llfig@p{\i}\fi\fi}%
  \fi
  \ifnum\llfig@k>1
    \foreach \i in {1,...,\the\llfig@k} {\foreach \j in {1,...,\the\llfig@k} {%
      \ifnum\i<\j
        \pgfmathparse{(\llfig@col{\i}{y}*\llfig@col{\j}{z} - \llfig@col{\i}{z}*\llfig@col{\j}{y})^2
          + (\llfig@col{\i}{z}*\llfig@col{\j}{x} - \llfig@col{\i}{x}*\llfig@col{\j}{z})^2
          + (\llfig@col{\i}{x}*\llfig@col{\j}{y} - \llfig@col{\i}{y}*\llfig@col{\j}{x})^2 > 0.0001}%
        \ifnum\llfig@rank<2 \ifnum\pgfmathresult=1
          \gdef\llfig@rank{2}\xdef\llfig@p{\i}\xdef\llfig@q{\j}\fi\fi
      \fi}}%
  \fi
  \ifnum\llfig@k>2
    \foreach \i in {1,...,\the\llfig@k} {\foreach \j in {1,...,\the\llfig@k} {%
      \foreach \l in {1,...,\the\llfig@k} {%
        \ifnum\i<\j \ifnum\j<\l
          \pgfmathparse{abs(
              \llfig@col{\i}{x}*(\llfig@col{\j}{y}*\llfig@col{\l}{z} - \llfig@col{\j}{z}*\llfig@col{\l}{y})
            + \llfig@col{\i}{y}*(\llfig@col{\j}{z}*\llfig@col{\l}{x} - \llfig@col{\j}{x}*\llfig@col{\l}{z})
            + \llfig@col{\i}{z}*(\llfig@col{\j}{x}*\llfig@col{\l}{y} - \llfig@col{\j}{y}*\llfig@col{\l}{x}))
            > 0.0001}%
          \ifnum\pgfmathresult=1 \gdef\llfig@rank{3}\fi
        \fi\fi}}}%
  \fi}
% The nonzero columns as labelled arrows. When the span is a line (in R^2
% or R^3), tip labels would sit on the line, so they go beside the arrows.
\def\llfig@drawcolumns{%
  \def\llfig@lp{tip}%
  \ifnum\llfig@rank=1 \ifnum\llfig@n>1 \def\llfig@lp{left}\fi\fi
  \ifnum\llfig@k>0
    \foreach \i in {1,...,\the\llfig@k} {%
      \pgfmathparse{(\llfig@col{\i}{x})^2 + (\llfig@col{\i}{y})^2
                    + (\llfig@col{\i}{z})^2 > 0.0001}%
      \ifnum\pgfmathresult=1
      \ifnum\llfig@n=3
        \pic {vector={to={\llfig@col{\i}{x},\llfig@col{\i}{y},\llfig@col{\i}{z}},
                      label=$\mathbf{\llfig@vname}_{\i}$, label pos=\llfig@lp}};
      \else\ifnum\llfig@n=2
        \pic {vector={to={\llfig@col{\i}{x},\llfig@col{\i}{y}},
                      label=$\mathbf{\llfig@vname}_{\i}$, label pos=\llfig@lp}};
      \else % R^1: arrows stacked just above the line
        \pic {vector={from={0,0.35*\i}, to={\llfig@col{\i}{x},0.35*\i},
                      label=$\mathbf{\llfig@vname}_{\i}$}};
      \fi\fi
      \fi}%
  \fi}
\newcommand{\SubspaceFig}[2][]{%
  \begingroup
  \pgfkeys{/llfig, reset, scale=1, color=blue, window=2.5, ambient=,
           vector name=v, dim caption=true, #1}%
  \global\llfig@k=0 \gdef\llfig@n{2}%
  \foreach \llfig@v in {#2} {%
    \llfig@ifnonempty\llfig@v{%
      \global\advance\llfig@k by 1
      \expandafter\llfig@readvec\llfig@v,,,\relax}}%
  \llfig@ifnonempty\llfig@ambient{\let\llfig@n\llfig@ambient}%
  \llfig@findrank
  % caption: dim W = r
  \ifllfig@dimcap % unless label below was given; the name has its own $s
    \if\relax\detokenize\expandafter{\llfig@below}\relax
      \if\relax\detokenize\expandafter{\llfig@name}\relax
        \def\llfig@below{$\dim = \llfig@rank$}%
      \else
        \def\llfig@below{$\dim$\ \llfig@name\ $= \llfig@rank$}%
      \fi
    \fi
  \fi
  % copies for passing into inner pics, whose own key reset would
  % otherwise redefine these in terms of themselves
  \let\llfig@sname\llfig@name \let\llfig@scolor\llfig@color
  \let\llfig@sns\llfig@ns \let\llfig@snt\llfig@nt
  \edef\llfig@L{\llfig@window}%
  \ifnum\llfig@n=1 % ---------------- R^1
    \begin{tikzpicture}[baseline, llfig scale=\llfig@scale, >=stealth]
      \begin{scope}[local bounding box=llfig@whole]
        \pgfmathtruncatemacro\llfig@Li{\llfig@L}%
        \pic {number line={min=-\llfig@Li, max=\llfig@Li, x1 label=$\mathbb{R}$}};
        \ifnum\llfig@rank=0
          \fill[color=\llfig@color] (0,0) circle (2.5pt);
        \else
          \draw[line width=3pt, color=\llfig@color, opacity=0.5]
            ({-\llfig@L-0.5},0) -- ({\llfig@L+0.5},0);
        \fi
        \llfig@drawcolumns
      \end{scope}
      \def\llfig@captionx{llfig@whole.center}%
      \llfig@captions[llfig@whole]
    \end{tikzpicture}%
  \else\ifnum\llfig@n=2 % ---------------- R^2
    \begin{tikzpicture}[baseline, llfig scale=\llfig@scale, >=stealth]
      \begin{scope}[local bounding box=llfig@whole]
        \ifnum\llfig@rank=2
          \fill[color=\llfig@color, opacity=0.15]
            (-\llfig@L,-\llfig@L) rectangle (\llfig@L,\llfig@L);
          \llfig@ifnonempty\llfig@name{%
            \node[anchor=north west, text=\llfig@color]
              at ({-\llfig@L},{\llfig@L}) {\llfig@name};}
        \fi
        \pic {axes 2d={xmin=-\llfig@L, xmax=\llfig@L,
                       ymin=-\llfig@L, ymax=\llfig@L}};
        \ifnum\llfig@rank=1 % the line s*d, s chosen to run just past the window
          \pgfmathsetmacro\llfig@s{(\llfig@L + 0.4)/max(abs(\llfig@col{\llfig@p}{x}),
                                                abs(\llfig@col{\llfig@p}{y}))}%
          \draw[very thick, color=\llfig@color, opacity=0.6]
            ({-\llfig@s*\llfig@col{\llfig@p}{x}},{-\llfig@s*\llfig@col{\llfig@p}{y}})
            -- ({\llfig@s*\llfig@col{\llfig@p}{x}},{\llfig@s*\llfig@col{\llfig@p}{y}})
            \llfig@ifnonempty\llfig@name{%
              node[pos=0.92, auto=right, text=\llfig@color] {\llfig@name}};
        \fi
        \ifnum\llfig@rank=0 \fill[color=\llfig@color] (0,0) circle (2.5pt); \fi
        \llfig@drawcolumns
      \end{scope}
      \def\llfig@captionx{llfig@whole.center}%
      \llfig@captions[llfig@whole]
    \end{tikzpicture}%
  \else % ---------------- R^3
    \begin{tikzpicture}[baseline, llfig view={\llfig@az}{\llfig@el},
                        llfig scale=\llfig@scale, >=stealth]
      \begin{scope}[local bounding box=llfig@whole]
        \pic {axes 3d={max={\llfig@L+0.6}}};
        \ifnum\llfig@rank=3
          \pic at (-\llfig@L,-\llfig@L,-\llfig@L) {parallelepiped={
            u={2*\llfig@L,0,0}, v={0,2*\llfig@L,0}, w={0,0,2*\llfig@L},
            vectors=false, fill=\llfig@scolor!40, opacity=0.25,
            style={draw=\llfig@scolor!60}}};
          \llfig@ifnonempty\llfig@name{%
            \node[text=\llfig@color, anchor=south west]
              at (\llfig@L,-\llfig@L,\llfig@L) {\llfig@name};}
        \fi
        \ifnum\llfig@rank=2 % square patch from an orthonormal basis of the span
          \pgfmathsetmacro\llfig@na{sqrt((\llfig@col{\llfig@p}{x})^2
            + (\llfig@col{\llfig@p}{y})^2 + (\llfig@col{\llfig@p}{z})^2)}%
          \pgfmathsetmacro\llfig@ex{\llfig@col{\llfig@p}{x}/\llfig@na}%
          \pgfmathsetmacro\llfig@ey{\llfig@col{\llfig@p}{y}/\llfig@na}%
          \pgfmathsetmacro\llfig@ez{\llfig@col{\llfig@p}{z}/\llfig@na}%
          \pgfmathsetmacro\llfig@dot{\llfig@col{\llfig@q}{x}*\llfig@ex
            + \llfig@col{\llfig@q}{y}*\llfig@ey + \llfig@col{\llfig@q}{z}*\llfig@ez}%
          \pgfmathsetmacro\llfig@fx{\llfig@col{\llfig@q}{x} - \llfig@dot*\llfig@ex}%
          \pgfmathsetmacro\llfig@fy{\llfig@col{\llfig@q}{y} - \llfig@dot*\llfig@ey}%
          \pgfmathsetmacro\llfig@fz{\llfig@col{\llfig@q}{z} - \llfig@dot*\llfig@ez}%
          \pgfmathsetmacro\llfig@nf{sqrt((\llfig@fx)^2 + (\llfig@fy)^2 + (\llfig@fz)^2)}%
          \pic {plane={u={\llfig@L*\llfig@ex, \llfig@L*\llfig@ey, \llfig@L*\llfig@ez},
                       v={\llfig@L*\llfig@fx/\llfig@nf, \llfig@L*\llfig@fy/\llfig@nf,
                          \llfig@L*\llfig@fz/\llfig@nf},
                       fill=\llfig@scolor!30, name=\llfig@sname,
                       name pos={\llfig@sns}{\llfig@snt}}};
        \fi
        \ifnum\llfig@rank=1
          \pgfmathsetmacro\llfig@s{(\llfig@L + 0.4)/max(abs(\llfig@col{\llfig@p}{x}),
            abs(\llfig@col{\llfig@p}{y}), abs(\llfig@col{\llfig@p}{z}))}%
          \draw[very thick, color=\llfig@color, opacity=0.6]
            ({-\llfig@s*\llfig@col{\llfig@p}{x}},{-\llfig@s*\llfig@col{\llfig@p}{y}},
             {-\llfig@s*\llfig@col{\llfig@p}{z}})
            -- ({\llfig@s*\llfig@col{\llfig@p}{x}},{\llfig@s*\llfig@col{\llfig@p}{y}},
                {\llfig@s*\llfig@col{\llfig@p}{z}})
            \llfig@ifnonempty\llfig@name{%
              node[pos=0.95, auto=right, text=\llfig@color] {\llfig@name}};
        \fi
        \ifnum\llfig@rank=0 \fill[color=\llfig@color] (0,0,0) circle (2.5pt); \fi
        \llfig@drawcolumns
      \end{scope}
      \def\llfig@captionx{llfig@whole.center}%
      \llfig@captions[llfig@whole]
    \end{tikzpicture}%
  \fi\fi
  \endgroup}

% ---------------- Schematics of Nul A, Col A and the fundamental subspaces
\newcommand{\KernelRangeFig}[1][]{%
  \begingroup
  \pgfkeys{/llfig, reset, scale=1,
    domain=$\mathbb{R}^n$, codomain=$\mathbb{R}^m$,
    kernel={$\operatorname{Nul} A$}, range={$\operatorname{Col} A$},
    map label={$\mathbf{x} \mapsto A\mathbf{x}$},
    point=$\mathbf{x}$, image=$A\mathbf{x}$, #1}%
  \begin{tikzpicture}[baseline, llfig scale=\llfig@scale, >=stealth]
    % the two spaces
    \filldraw[fill=gray!10, draw=gray!70, smooth cycle] plot coordinates
      {(-1.9,0.2) (-1.2,1.4) (0.4,1.6) (1.7,0.8) (1.6,-0.9) (0.2,-1.6)
       (-1.4,-1.1)};
    \filldraw[fill=gray!10, draw=gray!70, smooth cycle, xshift=6cm] plot
      coordinates {(-1.7,0.4) (-0.8,1.5) (0.9,1.3) (1.8,0.2) (1.2,-1.2)
       (-0.4,-1.5) (-1.5,-0.8)};
    \node[anchor=south east] at (-1.3,1.35) {\llfig@domain};
    \node[anchor=south west] at (6.8,1.3) {\llfig@codomain};
    % Nul A and Col A, each containing its 0
    \filldraw[fill=blue!20, draw=blue!60, smooth cycle] plot coordinates
      {(-0.2,-0.1) (0.5,0.4) (1.1,0.1) (1.0,-0.7) (0.2,-0.9)};
    \node[blue!70!black, anchor=north] at (0.5,-0.95) {\llfig@kernel};
    \filldraw[fill=red!15, draw=red!60, smooth cycle] plot coordinates
      {(5.0,-0.3) (5.4,0.6) (6.4,0.9) (7.1,0.3) (6.7,-0.6) (5.7,-0.8)};
    \node[red!70!black, anchor=north] at (6.1,-0.85) {\llfig@range};
    \fill (0.35,-0.35) circle (1.5pt) node[left] {$\mathbf{0}$};
    \fill (5.6,-0.2) circle (1.5pt) node[above] {$\mathbf{0}$};
    % x |-> Ax, and a vector of Nul A |-> 0
    \fill (-1.0,0.6) circle (1.5pt) node[below] {\llfig@ptlabel};
    \fill (6.4,0.4) circle (1.5pt) node[right] {\llfig@imlabel};
    \draw[->, thick, shorten <=2pt, shorten >=2pt] (-1.0,0.6)
      to[out=40, in=150] node[above] {\llfig@maplabel} (6.4,0.4);
    \fill (0.8,-0.4) circle (1.5pt);
    \draw[->, dashed, shorten <=2pt, shorten >=2pt] (0.8,-0.4)
      to[out=-20, in=-150] (5.6,-0.2);
  \end{tikzpicture}%
  \endgroup}
% One side of the fundamental-subspaces picture: a plane (#1) with the
% perpendicular line (#2) through 0; #3, #4 are optional dimension notes.
\newcommand{\llfig@perpside}[4]{%
  \begin{tikzpicture}[baseline, llfig view={30}{25}, llfig scale=\llfig@scale,
                      >=stealth]
    % the part of the line below the plane, seen through it
    \draw[very thick, red!80] (0,0,-1.1) -- (0,0,0);
    \pic {plane={u={1,0,0}, v={0,1,0}, extent=1.3, fill=blue!20,
                 opacity=0.7}};
    \draw[very thick, red!80] (0,0,0) -- (0,0,1.6);
    \draw (0.25,0,0) -- (0.25,0,0.25) -- (0,0,0.25); % right angle
    \fill (0,0,0) circle (1.5pt) node[below left=-1pt] {$\mathbf{0}$};
    % names, with the optional dimension notes as separate nodes
    \node[text=blue!60!black, inner sep=1pt] (llfig@pn) at (-0.6,0.8,0) {#1};
    \node[anchor=north, text=blue!60!black, font=\footnotesize, inner sep=1pt]
      at (llfig@pn.south) {#3};
    \node[anchor=south, text=red!80!black] (llfig@ln) at (0,0,1.6) {#2};
    \node[anchor=south, text=red!80!black, font=\footnotesize, inner sep=1pt]
      at (llfig@ln.north) {#4};
  \end{tikzpicture}}
\newcommand{\FundamentalSubspacesFig}[1][]{%
  \begingroup
  \pgfkeys{/llfig, reset, scale=1,
    domain=$\mathbb{R}^n$, codomain=$\mathbb{R}^m$,
    row space={$\operatorname{Row} A$}, kernel={$\operatorname{Nul} A$},
    range={$\operatorname{Col} A$}, left null={$\operatorname{Nul} A^T$},
    map label=$A$, dims=false, rank=r, #1}%
  \ifllfig@dims
    \def\llfig@da{$\dim = \llfig@rk$}\def\llfig@db{$\dim = n - \llfig@rk$}%
    \def\llfig@dc{$\dim = \llfig@rk$}\def\llfig@dd{$\dim = m - \llfig@rk$}%
  \else
    \let\llfig@da\empty\let\llfig@db\empty
    \let\llfig@dc\empty\let\llfig@dd\empty
  \fi
  \begin{tabular}{@{}c@{}}
    \llfig@perpside{\llfig@rowsp}{\llfig@kernel}{\llfig@da}{\llfig@db}\\[3pt]
    \llfig@domain
  \end{tabular}%
  \quad$\xrightarrow{\;\text{\llfig@maplabel}\;}$\quad
  \begin{tabular}{@{}c@{}}
    \llfig@perpside{\llfig@range}{\llfig@leftnul}{\llfig@dc}{\llfig@dd}\\[3pt]
    \llfig@codomain
  \end{tabular}%
  \endgroup}

% ---------------- The actual Nul A and Col A of a small matrix
% Rows are read with \llfig@readvec (so \llfig@k = m rows, \llfig@n = n
% entries per row) and their rank found with \llfig@findrank.
% \llfig@addvec{x}{y}{z}: append a vector, scaled to length 1.5, to the
% Nul A basis list (only the first \llfig@nA entries are kept)
\def\llfig@addvec#1#2#3{%
  \pgfmathsetmacro\llfig@len{sqrt((#1)^2 + (#2)^2 + (#3)^2)}%
  \pgfmathsetmacro\llfig@nvx{1.5*(#1)/\llfig@len}%
  \pgfmathsetmacro\llfig@nvy{1.5*(#2)/\llfig@len}%
  \pgfmathsetmacro\llfig@nvz{1.5*(#3)/\llfig@len}%
  \ifcase\llfig@nA\or
    \xdef\llfig@nullbasis{\llfig@nullbasis{\llfig@nvx},}%
  \or
    \xdef\llfig@nullbasis{\llfig@nullbasis{\llfig@nvx,\llfig@nvy},}%
  \else
    \xdef\llfig@nullbasis{\llfig@nullbasis{\llfig@nvx,\llfig@nvy,\llfig@nvz},}%
  \fi}
\newcommand{\NulColFig}[2][]{%
  \begingroup
  \pgfkeys{/llfig, reset, scale=1, matrix name=A, window=2, #1}%
  \global\llfig@k=0 \gdef\llfig@n{2}%
  \foreach \llfig@v in {#2} {%
    \llfig@ifnonempty\llfig@v{%
      \global\advance\llfig@k by 1
      \expandafter\llfig@readvec\llfig@v,,,\relax}}%
  \edef\llfig@nA{\llfig@n}\edef\llfig@mA{\the\llfig@k}%
  \llfig@findrank % rank of the rows, with independent rows p (and q)
  % columns of A, as a list for \SubspaceFig, and the largest entry
  \gdef\llfig@cols{}\gdef\llfig@big{1}%
  \foreach \j in {x,y,z} {%
    \ifnum\if\j x1\else\if\j y2\else3\fi\fi>\llfig@nA\relax\else
      \gdef\llfig@colj{}%
      \foreach \i in {1,...,\llfig@mA} {%
        \xdef\llfig@colj{\llfig@colj\ifnum\i>1 ,\fi\llfig@col{\i}{\j}}%
        \pgfmathsetmacro\llfig@tmp{max(\llfig@big, abs(\llfig@col{\i}{\j}))}%
        \global\let\llfig@big\llfig@tmp}%
      \xdef\llfig@cols{\llfig@cols{\llfig@colj},}%
    \fi}%
  % Nul A = (Row A)^perp, by the row rank
  \gdef\llfig@nullbasis{}%
  \edef\llfig@px{\llfig@col{\llfig@p}{x}}\edef\llfig@py{\llfig@col{\llfig@p}{y}}%
  \edef\llfig@pz{\llfig@col{\llfig@p}{z}}%
  \ifcase\llfig@nA\or % n = 1
    \ifnum\llfig@rank=0 \llfig@addvec{1}{0}{0}\fi
  \or % n = 2
    \ifcase\llfig@rank \llfig@addvec{1}{0}{0}\llfig@addvec{0}{1}{0}%
    \or \llfig@addvec{-(\llfig@py)}{\llfig@px}{0}\fi
  \else % n = 3
    \ifcase\llfig@rank
      \llfig@addvec{1}{0}{0}\llfig@addvec{0}{1}{0}\llfig@addvec{0}{0}{1}%
    \or % two vectors perpendicular to the row p: p x e_k, then p x (p x e_k),
        % with e_k for p's smallest entry so that p x e_k is not 0
      \pgfmathparse{and(abs(\llfig@px) <= abs(\llfig@py), abs(\llfig@px) <= abs(\llfig@pz))
        ? 1 : (abs(\llfig@py) <= abs(\llfig@pz) ? 2 : 3)}%
      \pgfmathtruncatemacro\llfig@kk{\pgfmathresult}%
      \ifcase\llfig@kk\or
        \edef\llfig@bx{0}\edef\llfig@by{\llfig@pz}\edef\llfig@bz{-(\llfig@py)}% p x e1
      \or
        \edef\llfig@bx{-(\llfig@pz)}\edef\llfig@by{0}\edef\llfig@bz{\llfig@px}% p x e2
      \else
        \edef\llfig@bx{\llfig@py}\edef\llfig@by{-(\llfig@px)}\edef\llfig@bz{0}% p x e3
      \fi
      \llfig@addvec{\llfig@bx}{\llfig@by}{\llfig@bz}%
      \llfig@addvec{\llfig@py*(\llfig@bz) - \llfig@pz*(\llfig@by)}%
                   {\llfig@pz*(\llfig@bx) - \llfig@px*(\llfig@bz)}%
                   {\llfig@px*(\llfig@by) - \llfig@py*(\llfig@bx)}%
    \or % p x q for two independent rows
      \edef\llfig@qx{\llfig@col{\llfig@q}{x}}\edef\llfig@qy{\llfig@col{\llfig@q}{y}}%
      \edef\llfig@qz{\llfig@col{\llfig@q}{z}}%
      \llfig@addvec{\llfig@py*(\llfig@qz) - \llfig@pz*(\llfig@qy)}%
                   {\llfig@pz*(\llfig@qx) - \llfig@px*(\llfig@qz)}%
                   {\llfig@px*(\llfig@qy) - \llfig@py*(\llfig@qx)}%
    \fi
  \fi
  % Col A shows the true columns; if they are long, its picture is scaled
  % down so that both halves come out about the same size
  \pgfmathsetmacro\llfig@colwin{max(\llfig@window, \llfig@big + 0.3)}%
  \pgfmathsetmacro\llfig@colscale{\llfig@scale*\llfig@window/\llfig@colwin}%
  \edef\llfig@tmp{%
    \noexpand\SubspaceFig[\unexpanded{#1}, ambient=\llfig@nA,
      name={$\noexpand\operatorname{Nul}\noexpand\,\llfig@mname$},
      vector name=u, color=blue]{\llfig@nullbasis}%
    \noexpand\quad$\noexpand\xrightarrow{\noexpand\;\llfig@mname\noexpand\;}$%
    \noexpand\quad
    \noexpand\SubspaceFig[\unexpanded{#1}, ambient=\llfig@mA,
      name={$\noexpand\operatorname{Col}\noexpand\,\llfig@mname$},
      vector name=a, color=red, window=\llfig@colwin,
      scale=\llfig@colscale]{\llfig@cols}}%
  \llfig@tmp
  \endgroup}
\makeatother
